\caption{Capacity--competence dissociation in M110. The markers and vertical ranges show the recorded minimum--maximum \reach sizes over six worlds. Labels show row-5 competence. Capacity increases from \mzero to \mtwo while row-5 success falls from 6/6 to 0/6 after the first acquired machinery generation. The vertical metric and the competence annotation are intentionally not combined into a single score.}
\label{fig:capacity}
\end{figure}

This distinction matters for self-improving systems. Enlarging a search space, action repertoire, or modification language can be a real capacity gain while worsening the policy that chooses among those possibilities. The result therefore argues against using a monotone capacity metric as a surrogate for realized competence under distribution shift.

\subsection{Alternative explanations and controls}
M110 includes several controls aimed at narrower explanations:
\begin{itemize}
  \item a deeper node bound of 13 leaves the fresh predecessor failures at rows 7 and 3 unchanged, so those are not shallow search failures;
  \item monotone-closure and visible-function certificates establish structural boundaries of the held spaces;
  \item fixed-point checks run at bounds 7, 9, 11, and 13 on every canonical world;
  \item removing generation two returns byte-exactly to \mone and loses row 3;
  \item removing generation one returns byte-exactly to \mzero, loses row 7, and regains row 5;
  \item a semantic mutation of generation two loses row 3;
  \item a built-but-unregistered rule leaves the state without that acquisition and does not obtain the claimed behavior;
  \item invalid state identity fails closed;
  \item all compared arms share world, demand, adapter, and budget bytes and differ in the rule cascade;
  \item host-widened controls show that the host can directly widen the named components, clarifying that the claim is about attribution and state-held machinery rather than a capability the host itself lacks.
\end{itemize}
All P1--P24 compute true with replay equality across 186 isolated processes and zero external calls. The consumer family remains project-authored and was deliberately chosen to realize the stress row. M110 is therefore a designed causal stress test, not independent-domain transfer.

\section{M111: self-directed diagnosis under a scarce experiment budget}
\label{sec:m111}
M110 shows a failure of extrapolation: the inherited feature row can drive a confident but wrong attribution outside the producer geometry. M111 asks a more precise successor question. Can the lineage use its own accumulated record to identify observations for which the feature vocabulary does not determine the answer, and can it spend a limited experiment only there?

\subsection{An information bound, not a noisy prediction}
